\documentclass{article}
\usepackage{a4wide}
\usepackage{geometry}
\geometry{
   includeheadfoot,
   margin=2.54cm
}
\usepackage[utf8]{inputenc}
\usepackage{hyperref}
\usepackage[spanish]{babel}

\title{README para repositorio de Matlab}
\author{Diego Fernández de la Peña \\
Tutor: Uwe Brauer
}
\date{\today}

\makeatletter
\usepackage{paralist}
\usepackage[colorinlistoftodos, textwidth=3cm, shadow]{todonotes}


\definecolor{GreenYellow}{rgb}{0.678431,1.000000,0.184314}
\newcounter{todocounter}
\newcommand{\todonum}[2][]
{\stepcounter{todocounter}\todo[#1]{\thetodocounter: #2}}

\newcounter{ubcomment}

 
\newcommand{\ubcomment}[2][]{%
\refstepcounter{ubcomment}%
{%
\todo[fancyline,linecolor=black,backgroundcolor={red!40!},size=\footnotesize]{%
\textbf{Fixme: U.B [\uppercase{#1}\theubcomment]:}~#2}%
}}



\newcommand{\ubresolved}[2][]{%
\refstepcounter{ubcomment}%
{%
\todo[linecolor=black,backgroundcolor={red!40!},size=\footnotesize]{%
\textbf{Resolved: UB [\uppercase{#1}\theubcomment]:}~#2}%
}}

\newcommand{\ubcommentinline}[2][]{%
\refstepcounter{ubcomment}%
{%
\todo[linecolor=black,inline,backgroundcolor={red!40!},size=\footnotesize]{%
\textbf{Fixme: U.B. [\uppercase{#1}\theubcomment]:}~#2}%
}}


\newcommand{\ubresolvedinline}[2][]{%
\refstepcounter{ubcomment}%
{%
\todo[linecolor=black,inline,backgroundcolor={red!40!},size=\footnotesize]{%
\textbf{Resolved: UB [\uppercase{#1}\theubcomment]:}~#2}%
}}


\newcommand{\ubcommentmultiline}[2]{%
\refstepcounter{ubcomment}%
{%
\todo[linecolor=black,inline,caption={\textbf{{Fixme: U.B.}
    [\theubcomment]:} #1} ,backgroundcolor={red!40!},size=\footnotesize]{%
\textbf{Fixme: U.B. [\theubcomment]:}~#2}%
}}


\newcounter{dfcomment}
\newcommand{\dfcomment}[2][]{%
\refstepcounter{dfcomment}%
{%
\todo[fancyline,color={purple!50!}]{%
\textbf{Fixme: DF [\uppercase{#1}\thedfcomment]:}~#2}%
}}

\newcommand{\dfcommentinline}[2][]{%
\refstepcounter{dfcomment}%
{%
%\todo[linecolor=black,inline,caption={},backgroundcolor={blue!90!},size=\footnotesize]{%  
\todo[fancyline,linecolor=black,inline,backgroundcolor={yellow!80!},size=\footnotesize]{%  
\textbf{Fixme: DF [\uppercase{#1}\thedfcomment]:}~#2}%
}}




\newcommand{\dfcommentmultiline}[2]{%
\refstepcounter{dfcomment}%
{%
\todo[inline,caption={\textbf{{Fixme: DF}
    [\thedfcomment]: #1 }},color={blue!40!}]{%
\textbf{Fixme: DF [\thedfcomment]:}~#2}%
}}



\makeatother

\usepackage{marginnote}
\makeatletter
\renewcommand{\@todonotes@drawMarginNoteWithLine}{%
\begin{tikzpicture}[remember picture, overlay, baseline=-0.75ex]%
    \node [coordinate] (inText) {};%
\end{tikzpicture}%
\marginnote[{
    \@todonotes@drawMarginNote%
    \@todonotes@drawLineToLeftMargin%
}]{% Draw note in right margin
    \@todonotes@drawMarginNote%
    \@todonotes@drawLineToRightMargin%
}%
}
\makeatother


\begin{document}

\maketitle
\section{Cómo usar el repositorio}

Este repositorio contiene implementaciones de distintos métodos numéricos 
para resolver la ecuación de Lane-Emden. No existe un programa principal, 
cada método se ejecuta de forma separada, salvo en el caso de los métodos 
SLM y QLM, que se prueban conjuntamente.

\subsection{Métodos SLM y QLM}

Los métodos SLM y QLM se prueban mediante el script:
\begin{verbatim}
>> laneemdenprueba
\end{verbatim}
Este archivo \texttt{laneemdenprueba.m} está diseñado para:
\begin{itemize}
  \item Ejecutar el método \textbf{SLM}.
  \item Ejecutar el método \textbf{QLM}.
\end{itemize}
No incluye los métodos \texttt{ode45} ni \texttt{RK4}. 
Su función es únicamente validar y analizar el comportamiento de 
los métodos iterativos. Cuando se ejecuta este script se obtiene 
como resultado una tabla con la solución aproximada para los 
distintos valores del índice politrópico en cada iteración.
Depende del método que se quiera probar, se deberá de introducir
\texttt{laneemdenSLM.m} o \texttt{laneemdenNEWTONKANTOROVICH.m}
en la línea correcta del script.

\subsection{Métodos clásicos \texttt{ode45} y RK4}

Los métodos clásicos se ejecutan de forma independiente, llamando 
directamente a sus respectivos archivos:
\begin{itemize}
  \item \texttt{laneemdenode45.m}: resuelve la ecuación de Lane-Emden usando 
  \texttt{ode45} de Matlab.
  \item \texttt{laneemdenRK4.m}: resuelve la ecuación de Lane-Emden mediante 
  el método de Runge-Kutta clásico de orden 4 con paso fijo.
\end{itemize}
Estos archivos no dependen de \texttt{laneemdenprueba.m} y pueden 
ejecutarse directamente desde la consola de Matlab proporcionando 
los parámetros adecuados (como el índice politrópico y el intervalo 
de integración).

\subsection{Ejecutar un método específico}

Si se quiere ejecutar un método en particular, hay que invocarlo directamente. 
Por ejemplo, con el comando de Matlab:
\begin{verbatim}
>> theta = laneemdenode45(n, epsilon, puntofinal);
\end{verbatim}
en el que se tiene que consultar el comentario inicial 
dentro de la función para ver la sintaxis exacta de sus parámetros.

\subsection{Visualización de resultados}

La visualización se realiza mediante el script:
\begin{verbatim}
>> graficas
\end{verbatim}
Al ejecutar \texttt{graficas.m}, el programa genera automáticamente 
\textbf{tres figuras}:
\begin{itemize}
  \item Una figura con los resultados obtenidos por los métodos \textbf{SLM y QLM}.
  \item Una figura con las soluciones obtenidas mediante \textbf{ode45} y \textbf{RK4}.
  \item Una figura con las soluciones analíticas de la ecuación de Lane-Emden.
\end{itemize}
Este archivo permite comparar visualmente el comportamiento de los distintos 
métodos numéricos y el de las soluciones exactas.
Un punto a tener muy en cuenta es que para poder visualizar la gráfica
de las funciones \textbf{ode45} y \textbf{RK4} es necesario leer el
comentario final de ambas funciones, ya que si no se tiene en cuenta
solamente representará un punto, el primer cero de la ecuación. Esto
es debido a que ambas funciones están programadas para que devuelvan
el punto donde se consigue el primer cero, por ello, para una buena 
representación es necesario el vector entero, y para ello, se necesita 
leer el comentario al que antes se ha hecho referencia para ejecutarlo
correctamente.

\section{Qué puede esperar el usuario}

Al utilizar este repositorio, el usuario puede:
\begin{itemize}
  \item Resolver numéricamente la ecuación de Lane-Emden mediante 
  métodos iterativos y métodos clásicos.
  \item Comparar gráfica y analíticamente SLM y QLM.
  \item Comparar gráficamente \texttt{ode45} y RK4.
  \item Contrastar las soluciones analíticas con las soluciones numéricas.
\end{itemize}

\end{document}

